\section{Open issues}
\label{sec:open}
Ordered by importance for the next step.
\begin{enumerate}
\item [decision] Matching rule for the final correction (cell$\pm$1 vs cell) and mixing classes.
\item Improve the low-$p_T$ projection (energy loss, or outer helix not available in picoDst).
\item Translate $P_{\rm ME}$ into a random-match correction of the TOF-matched yield and \mtwo{} background
  (needs the time of the random hit, not stored in picoDst).
\item Line shapes: no fit model is statistically acceptable below 2\,GeV/$c$ on the large sample (Sec.~\ref{sec:bigfit}); try finer
  momentum bins or a resolution-scaled variable, and understand the continuum left at 0.2--0.4\,GeV/$c$.
\item Track-correlated late hits (Sec.~\ref{sec:lowp}): a charge-symmetric SE excess of about \lowpTwoADmin--\lowpTwoADmax{} per
  $10^4$ tracks in $1.6<m^2<2.8$\,GeV$^2$/$c^4$ at 0.2--0.4\,GeV/$c$, independent of centrality, plus a positive-only part at
  0.2--0.3\,GeV/$c$. The full simulation attributes it mainly to secondaries from the BTOF trays (Sec.~\ref{sec:sim}); the
  positive-only part is not reproduced.
\item Simulation (Sec.~\ref{sec:sim}): ToT and the slow BTOF simulator were tested (Sec.~\ref{sec:simvar}); a ToT raises the
  0.4--0.6\,GeV/$c$ correlated and unrelated rates toward the data but none of the responses removes the 3--4$\times$ over-prediction of
  correlated matches in the K--p valley and above the deuteron at 0.2--0.4\,GeV/$c$. Open: GEANT3 settings for low-energy secondaries
  in the trays (time and energy cuts, hadronic package), electronics dead time, multiplicity tuning. Decided: simulation as evidence
  only, not as a correction (Sec.~\ref{sec:simtmpl}); the component is estimated from the data (Sec.~\ref{sec:corrdd}).
\item Data-driven correlated component (Sec.~\ref{sec:corrdd}): systematics ($\kappa$, i.e.\ path length, $\pm10\%$; separate amplitude
  for positive tracks; shape under the peaks) not evaluated; check the BTOF readout window against the fitted cut-off $t_0$.
  Decided: quoted as a systematic of the random-match background, not subtracted; it exceeds the 10\% target only for protons at
  0.2--0.4\,GeV/$c$ (\cddpAllA\%).
\item Method 2 loses \ovbigSElose\% of the own matches in the overlay; quantify the effect on the random fraction.
\item MuDst statistics: 98 files of 14 runs restored (\fpath{filelists/mudst/restore_100.txt}, 319\,GB); one job per file.
  Run 20088015 needs a workaround for the trigger-data decoding errors.
\item Understand the slight $p_T$ dependence of $P_{\rm ME}$ seen in the 1M sample.
\item Process the three removed files (030, 033, 040) if the full 57-file sample is needed.
\item [decision] Confirm the trigger list and decide on VPD-vertex and pile-up cuts.
\item Scale beyond 1M events if needed (one HTCondor job per run, about 70k events in 35--40 min; avoid the qgp017 server).
\item Understand the reduced TPC track density near $\phi\approx-0.6$, $\eta>0$.
\end{enumerate}
